%% This paper is heavily inspired by the work of Maximilian Bachl, Joachim Fabini and Tanja Zseby, as presented in their paper "A flow-based IDS using Machine Learning in eBPF". The original paper can be found at: http://arxiv.org/abs/2102.09980 and their .tex file on their github https://github.com/CN-TU/machine-learning-in-ebpf
%%
%% 
%%
%%
%% This is file `sample-sigconf.tex',
%% generated with the docstrip utility.
%%
%% The original source files were:
%%
%% samples.dtx  (with options: `sigconf')
%%
%% IMPORTANT NOTICE:
%%
%% For the copyright see the source file.
%%
%% Any modified versions of this file must be renamed
%% with new filenames distinct from sample-sigconf.tex.
%%
%% For distribution of the original source see the terms
%% for copying and modification in the file samples.dtx.
%%
%% This generated file may be distributed as long as the
%% original source files, as listed above, are part of the
%% same distribution. (The sources need not necessarily be
%% in the same archive or directory.)
%%
%% The first command in your LaTeX source must be the \documentclass command.
\documentclass[conference]{IEEEtran}

\usepackage[utf8]{inputenc}
\usepackage{multirow}
\usepackage[inline]{enumitem}
\usepackage{xcolor}
\usepackage{booktabs}
\usepackage{tabularx}
\usepackage[pdfauthor={Sacha Corporeau}]{hyperref}
\usepackage{amsmath}
\usepackage{amssymb}
\usepackage{graphicx}
\usepackage[numbers]{natbib}
\usepackage{subfig}
\usepackage{tikz}
\usepackage{float}
\usepackage[OT1]{fontenc}

\usepackage[nomain, toc, acronym]{glossaries}
\glsdisablehyper


\begin{document}

\title{ML for Intrusion Detection System in smartNICs using eBPF XDP}

\author{
    \IEEEauthorblockN{Sacha Corporeau}
    \IEEEauthorblockA{Student at Télécom SudParis \& Intern at Loria\\
    sacha.corporeau@loria.fr}
}
%\date{\today}

\maketitle

\newacronym{ml}{ML}{Machine Learning}
\newacronym{ids}{IDS}{Intrusion Detection System}
\newacronym{xdp}{XDP}{eXpress Data Paths}
\newacronym{ebpf}{eBPF}{extended BPF}
\newacronym{snic}{sNIC}{smart Network Interface Card}
\newacronym{mlp}{MLP}{Multi-Layer Perceptron}
\newacronym{nn}{NN}{Neural Network}
\newacronym{tc}{TC}{Traffic Control}
\newacronym{dt}{DT}{Decision Tree}
\newacronym{afxdp}{AF\_XDP}{Address Family Express Data Path}
\newacronym{gro}{GRO}{Generic Receive Offload}

\begin{abstract}
    Throughout the last 10 years, we have seen many papers exploring \gls{ebpf} in order to implement various things. One of the studied fields is \gls{ids} which uses \gls{ebpf}'s \gls{xdp} in order to implement packet filtration inside the linux kernel. In parallel, the growing use of \gls{snic} with \gls{xdp} offloading allows to run bigger \gls{ebpf} scripts without impacting the host machine's CPU. In this paper, we purpose different implementations for many \gls{ml} techniques applied to \gls{ids} with \gls{xdp} scripts. We analyzed all kind of performances for each \gls{ml} technique that we choose to implement with different \gls{ebpf} techniques : \gls{xdp} (both generic and driver) and \gls{tc}. The goal is to proove that the efficient way is to use \gls{xdp} at the driver level and even offload it to \gls{snic}.

\end{abstract}


\section{State-of-the-art}

First, This works is heavily motivated by \cite{rajurkar_packets_2025} which highlight that \gls{ml} in \gls{ebpf} has some non-negligible constraints but has to be studied in-depth because of the huge benefits that \gls{snic}s offloading can afford. This paper reviewed many \gls{ebpf}-based \gls{ids} papers published between 2019 and 2025. It has identified 5 main challenges faced during many \gls{ids} implementation in \gls{ebpf} :
\begin{itemize}
    \item Verifier constraints and in-kernel programmability
    \item Cross-kernel portability
    \item \gls{ml} integration complexity
    \item Observability vs. privacy trade-offs
    \item Resource efficiency at scale
\end{itemize}

Zhang et. al \cite{zhang_real-time_2024} purposed an implementaion of a \gls{mlp} in-kernel making use of \gls{ebpf}'s tailing to allow the execution of a \gls{mlp} with three hidden layers, each of size 32. They focused on optimizing their method in order to archieve performances "comparable to the best \gls{ebpf}-based method".
They also implemented a way to hot-update the \gls{nn} parameters, in order to allow 100\% runtime of the \gls{ids} \gls{ebpf} program, thus preventing intrusion while rebooting the program.
This work is really interesting because it highlight that it is possible to implement some \gls{mlp} in-kernel, but lack of real benchmark with driver \gls{xdp} for example.

Then, \cite{hara_practicality_2024} purposed a more detailed implementation of \gls{mlp} and \gls{dt}, while comparing user-space and kernel-space behaviors. They used quitea range of \gls{ebpf} technologies \: from the use of \gls{afxdp} to native \gls{xdp}. More generally, they made a survey on all \gls{ml}-\gls{ebpf} based-papers, and all the optimizations around it (such as quantization, offloading on \gls{snic}s, \dots). They provided quite a lot of data to support their claims about But this work still show a lack of benchmark for driver \gls{xdp}.

Bachl et al. \cite{bachl_flow-based_2022} used their own python library to recreate virtual network interfaces virtio, that uses veth command from linux. Without adding more options, veth enable \gls{gro} by default and uses generic \gls{xdp}. Thus, both their kernel-space and user-space applications runs after the \gls{gro} fires in the linux kernel network stack, so they both use 65535 Bytes packets to run their program.



\begin{table*}
    \caption{Comparison of the different benchmarks of \gls{ebpf} made so far} \label{tab:benchmarks_comparison}
    \centering
    \begin{tabularx}{\textwidth}{XXXX}
        \toprule
        Paper & \gls{ebpf} mode & program behavior & Measured results\\
        \midrule
        \text{Kulkarni \cite{kulkarni_ebpf_2023,kulkarni_performance_2023}} & \gls{tc}, generic \gls{xdp}, native \gls{xdp} & Dummy & Throughput and FlameGraph \\
        \text{Zhang et al. \cite{zhang_real-time_2024}} & generic \gls{xdp}  & \gls{mlp} \gls{ids} & Classification metrics, execution times\\
        \text{Hara and Sasabe \cite{hara_practicality_2024}} & all : \gls{afxdp}, native \gls{xdp}, generic \gls{xdp} and \gls{tc} & \gls{nn}, \gls{dt}, simple extraction & \\
        \text{Bach et al. \cite{bachl_flow-based_2022}} & \gls{xdp} generic, user-space & Flow-based \gls{dt} & packets per second \\
        \bottomrule
    \end{tabularx}
\end{table*}

While reading about \gls{ml}-based \gls{ids} in \gls{ebpf}, we saw two different implementations for the \gls{ids}. Flow-based and Packet-based \gls{ids}. This is a critical choice when implementing \gls{ids} because Flow-based detection constraints the \gls{ebpf} program to use a significant part of the maps (often with a LRU Hash map) to register flow informations.

One of the issues of the previous work is that some performance measures were made in number of packets (execution time, packets processed per second, \dots), some other in throughput (Bytes per second). This would have been quite the same if the size of the packets processed were the same in all the different types of \gls{ebpf} hooks. But the issue here is the existence of \gls{gro}. \gls{gro} is a feature of the linux kernel introduced in XXXXX taht fires between the NIC and the \gls{tc} hook which reconstructs packet flows from the max size of the network interface (usually 1500 Bytes) to a new bigger packet size (usually 65535 Bytes) before forwading it to the user space (and the \gls{tc} hook in our case).
This might be not the only feature of the linux kernel or the NIC driver that can modify the behavior of \gls{ebpf} scripts. For this reason, it is interesting to realise some new benchmark of \gls{ebpf} technology for \gls{ml}-based \gls{ids} with all different kind of network interface parameters or \gls{ebpf} technologies.

\section{Objectives}
The aims of this work is to evaluate different performances of \gls{ml}-based \gls{ids} in \gls{ebpf} with all kind of possibilities. Different kernel hooks, memory optimisation, \dots

The second objective is to evaluate and study \gls{snic}s offloading of theses programs with some proof-of concept and performance evaluation.

\subsection{Different ML types}
TODO

\begin{itemize}
    \item model choice : DT, RF, CNN, MLP, \dots
    \item Flow-based or Packet-based \gls{ids}
\end{itemize}

\subsection{SmartNIC offloading}

TODO
Look at bibliography \& test on \gls{snic}.


\section{Results}

In order to benchmark connexion between pairs machines, we used two VMs in a proxmox cluster, both running Ubuntu 24.04 'Noble Numbat' with kernel 6.8.0-110-generic. The first server with ip 172.16.6.2 and second with ip 172.16.6.3.
They are virtually connected together with a theoratically 5Gbps interface, bu we will see what is the max throughput archieved in the results.

\subsection{First Benchmark}

First benchmark is a throughput Benchmark, analyzing max throughput archieved while running simple net\_listener scripts. In order to measure throughput, iperf3 showed to be an efficient and easy-to-use module in order to emulate max throughput \cite{bachl_flow-based_2022}. For each measurment, we ran five 30s iperf benchmarks with 10 parallel clients from the same server.

\begin{table}
    \caption{Results from the first benchmark on parsing times \& maximum throughput} \label{tab:results1}
    \centering
    \begin{tabular}{ccc}
        \toprule
        tracepoint & TC & XDP \\
        \midrule
        \bottomrule
    \end{tabular}
\end{table}

Then, we also wanted to measure execution time of code parts. Each ebpf script used is divided in 2 parts : parsing and classification. We carefuly noted kernel timestamps obtained with the ebpf helper function bpf\_ktime\_get\_ns. Now, we wanted to analyze the parsing time for the same network packets. Since the TC and XDP scirpts were completely able to run in parallel, the results shown under for theses two scritps were done with (almost) the exact same packets. Since running \gls{xdp} natively (not offloaded, but in the NIC driver) prevents the tracepoint netif\_receive\_skb to fire (I do not really know exactly why for now), we had to re\-run the same benchmark for the net\_listener.bpf script apart from the others. This can lead to more imprecision in the interpretation of the results, but because the tracepoint approach was provided only for a comparison (since it is not able to pass or drop packets with \gls{ebpf}), this is not a concern.

The packets sent during theses parsing times benchmark came from multiple sources :
\begin{itemize}
    \item ssh connection to my personal computer in order to launch and monitor the scripts
    \item another iperf3 throughput test similar to the one in the first benchmark
\end{itemize}

The results are shown in this log-histogram of parsing and total processing time per packet.


\bibliographystyle{ieeetr}
\bibliography{eBPF_on_smartNIC}


\end{document} 